\hypertarget{methodology}{%
\section{Methodology}\label{methodology}}

This work adopts a Design Science Research (DSR) approach (Hevner et
al., 2004): apart from the observation of the routing problem presented
in Section 2.1, the project required the design and testing of an
artefact, the context-aware agentic gateway, that would carry the whole
knowledge contribution. The build-evaluate DSR loop maps onto the
project stages: phase P5 and parts of P6-P7, namely simulation and PPO
offline training, and implementation of the infrastructure and
individual parts respectively correspond to the build cycle. Phase P8,
which involves comparing against five baselines, testing three scenario
variants, and a seven-run ablation (Section 7.1),
corresponds to the evaluate cycle. The product of this cycle is the
gateway, whereas an accompanying knowledge contribution is the six-item
list of features, benefits and capabilities set out in Section 2.3, each
of which is verified against a specific evaluation output rather than
only declared.

DSR only defines a broad criterion for what is considered valid
evidence. It does not set out the organisation of the weekly work tasks
inside each phase. Managing weekly tasks is handled by Scrum sprints
(Schwaber and Sutherland, 2020): each week\textquotesingle s sprint
backlog is allocated to one particular component development or one
evaluation run, and weekly log submission at Weeks 6, 9 and 12 is the
evidence of each sprint review. The two approaches are complementary: if
a checker wanted to look up a student\textquotesingle s learning outcome
2 achievement ("critically select and assess the most suitable
methodological processes"), the logbook entry for a given week can be
traced through the DSR stage the work fed into, without the two accounts
of how the work was done contradicting each other.

